% MIG instances on a seven-slice GPU (Lesson 9).  Seven GPU slices, each
% with its compute (GPU) and memory part; four applications own GPU
% instances of 3, 2, 1 and 1 slices.  Own allocation; the outline figure
% shows 1, 1, 2 and 3.  The memory side is simplified to one block per
% slice: the A100 has eight memory slices and a 3-slice instance gets
% four (3g.20gb), see the caption.
% Redraws outline figure m9l2_gpumulti (allocation changed).
% Usage: % MIG instances on a seven-slice GPU (Lesson 9).  Seven GPU slices, each
% with its compute (GPU) and memory part; four applications own GPU
% instances of 3, 2, 1 and 1 slices.  Own allocation; the outline figure
% shows 1, 1, 2 and 3.  The memory side is simplified to one block per
% slice: the A100 has eight memory slices and a 3-slice instance gets
% four (3g.20gb), see the caption.
% Redraws outline figure m9l2_gpumulti (allocation changed).
% Usage: % MIG instances on a seven-slice GPU (Lesson 9).  Seven GPU slices, each
% with its compute (GPU) and memory part; four applications own GPU
% instances of 3, 2, 1 and 1 slices.  Own allocation; the outline figure
% shows 1, 1, 2 and 3.  The memory side is simplified to one block per
% slice: the A100 has eight memory slices and a 3-slice instance gets
% four (3g.20gb), see the caption.
% Redraws outline figure m9l2_gpumulti (allocation changed).
% Usage: % MIG instances on a seven-slice GPU (Lesson 9).  Seven GPU slices, each
% with its compute (GPU) and memory part; four applications own GPU
% instances of 3, 2, 1 and 1 slices.  Own allocation; the outline figure
% shows 1, 1, 2 and 3.  The memory side is simplified to one block per
% slice: the A100 has eight memory slices and a 3-slice instance gets
% four (3g.20gb), see the caption.
% Redraws outline figure m9l2_gpumulti (allocation changed).
% Usage: \input{figures/l09-mig}
% Width: margin (50 mm); the drawing is about 48 mm wide.
\begin{tikzpicture}[x=1mm, y=1mm, font=\scriptsize\sffamily,
  g/.style={draw, fill=ln-accent!30, minimum width=6mm, minimum height=5mm, inner sep=0pt},
  m/.style={draw, fill=ln-tip!35, minimum width=6mm, minimum height=3.5mm, inner sep=0pt},
  app/.style={draw, fill=ln-box, minimum height=4.5mm, inner sep=0pt}]
  \foreach \i in {0,...,6} {
    \node[g] at (3+\i*7,10) {\the\numexpr\i+1\relax};
    \node[m] at (3+\i*7,5) {};
  }
  \node[app, minimum width=20mm] at (10,-1) {App 1};
  \node[app, minimum width=13mm] at (27.5,-1) {App 2};
  \node[app, minimum width=6mm] at (38,-1) {3};
  \node[app, minimum width=6mm] at (45,-1) {4};
  \node[anchor=west] at (-0.5,15.5) {\iflangja{GPU スライス（計算＋メモリ）}{GPU slices (compute + memory)}};
\end{tikzpicture}

% Width: margin (50 mm); the drawing is about 48 mm wide.
\begin{tikzpicture}[x=1mm, y=1mm, font=\scriptsize\sffamily,
  g/.style={draw, fill=ln-accent!30, minimum width=6mm, minimum height=5mm, inner sep=0pt},
  m/.style={draw, fill=ln-tip!35, minimum width=6mm, minimum height=3.5mm, inner sep=0pt},
  app/.style={draw, fill=ln-box, minimum height=4.5mm, inner sep=0pt}]
  \foreach \i in {0,...,6} {
    \node[g] at (3+\i*7,10) {\the\numexpr\i+1\relax};
    \node[m] at (3+\i*7,5) {};
  }
  \node[app, minimum width=20mm] at (10,-1) {App 1};
  \node[app, minimum width=13mm] at (27.5,-1) {App 2};
  \node[app, minimum width=6mm] at (38,-1) {3};
  \node[app, minimum width=6mm] at (45,-1) {4};
  \node[anchor=west] at (-0.5,15.5) {\iflangja{GPU スライス（計算＋メモリ）}{GPU slices (compute + memory)}};
\end{tikzpicture}

% Width: margin (50 mm); the drawing is about 48 mm wide.
\begin{tikzpicture}[x=1mm, y=1mm, font=\scriptsize\sffamily,
  g/.style={draw, fill=ln-accent!30, minimum width=6mm, minimum height=5mm, inner sep=0pt},
  m/.style={draw, fill=ln-tip!35, minimum width=6mm, minimum height=3.5mm, inner sep=0pt},
  app/.style={draw, fill=ln-box, minimum height=4.5mm, inner sep=0pt}]
  \foreach \i in {0,...,6} {
    \node[g] at (3+\i*7,10) {\the\numexpr\i+1\relax};
    \node[m] at (3+\i*7,5) {};
  }
  \node[app, minimum width=20mm] at (10,-1) {App 1};
  \node[app, minimum width=13mm] at (27.5,-1) {App 2};
  \node[app, minimum width=6mm] at (38,-1) {3};
  \node[app, minimum width=6mm] at (45,-1) {4};
  \node[anchor=west] at (-0.5,15.5) {\iflangja{GPU スライス（計算＋メモリ）}{GPU slices (compute + memory)}};
\end{tikzpicture}

% Width: margin (50 mm); the drawing is about 48 mm wide.
\begin{tikzpicture}[x=1mm, y=1mm, font=\scriptsize\sffamily,
  g/.style={draw, fill=ln-accent!30, minimum width=6mm, minimum height=5mm, inner sep=0pt},
  m/.style={draw, fill=ln-tip!35, minimum width=6mm, minimum height=3.5mm, inner sep=0pt},
  app/.style={draw, fill=ln-box, minimum height=4.5mm, inner sep=0pt}]
  \foreach \i in {0,...,6} {
    \node[g] at (3+\i*7,10) {\the\numexpr\i+1\relax};
    \node[m] at (3+\i*7,5) {};
  }
  \node[app, minimum width=20mm] at (10,-1) {App 1};
  \node[app, minimum width=13mm] at (27.5,-1) {App 2};
  \node[app, minimum width=6mm] at (38,-1) {3};
  \node[app, minimum width=6mm] at (45,-1) {4};
  \node[anchor=west] at (-0.5,15.5) {\iflangja{GPU スライス（計算＋メモリ）}{GPU slices (compute + memory)}};
\end{tikzpicture}
